\section{Crossing candidate scope and query budgets}
We separate the small slate menu from its frozen count restriction
on the same 24 development recordings and three seeds. With exact
counts we use the previously fixed number of edits in every query.
With relaxed counts the narrow arm can choose any length of the same
five causal prefix rankings, the original slate, or no edits. The
full bank can choose any legal path, at most one per agent, in either
count condition. All four use the original CV risk and support rule
and must beat CV and OLS4 on conservative full-grid utility.

\begin{center}
\small
\setlength{\tabcolsep}{3pt}
\begin{tabular}{@{}lrrrr@{}}
\toprule
Condition & Verified & Excluded & Unsupported & Open \\
\midrule
Narrow / exact & \CapacityNarrowExactStatus \\
Full / exact & \CapacityFullExactStatus \\
Narrow / relaxed & \CapacityNarrowRelaxedStatus \\
Full / relaxed & \CapacityFullRelaxedStatus \\
\bottomrule
\end{tabular}
\end{center}
Each row has \CapacityViews{} recording--seed views. Exclusions are
risk infeasibility or a negative conservative upper bound. When a
solve is unresolved, a verified witness from a nested condition is
carried forward only after rescoring it under the wider condition. With
frozen counts, expanding the menu changes the source-balanced
verified rate by \CapacityWidthFixed{} pp; relaxing counts within the
narrow menu changes it by \CapacityBudgetNarrow{} pp. At relaxed
counts, the full bank's change over the narrow menu is
\CapacityWidthRelaxed{} pp. Within the full bank the effect of
relaxing counts is \CapacityBudgetFull{} pp. Intervals are paired,
3,000-draw source bootstrap intervals for verified witness rates,
not intervals for unobserved capacity. The four comparisons describe
retrospective feasible sets, not deployable estimators or risk guarantees.
